\section{An angle minimum and its spectral gap}\label{sec:angle}

We now construct two spatial orbitals with the prescribed density.
The main point is to prove a spectral gap for their common formal
potential.  This will allow the next section to compare the resulting
determinant with arbitrary fermionic ensembles.

Write \(\Omega_+=\{x\in\R^3:z>0\}\).  For a positive continuous weight \(w\),
set
\[
 H^1(w)=
 \left\{\xi\in H^1_{\mathrm{loc}}(\R^3):
       \int_{\R^3}w\bigl(|\xi|^2+|\nabla\xi|^2\bigr)<\infty\right\},
 \qquad
 \mathcal E_w(\xi)=\frac12\int_{\R^3}w|\nabla\xi|^2.
\]
When \(w\) is integrable, use the unnormalized weighted variance
\[
 \Var_w(\xi)=
 \int_{\R^3}w|\xi|^2
 -\frac{\left|\int_{\R^3}w\xi\right|^2}{\int_{\R^3}w}.
\]
The corresponding spaces on \(\Omega_+\) are defined in the same way,
with local Sobolev regularity up to its boundary when a trace is specified.
On compact sets all these weighted norms are equivalent to the ordinary
Sobolev norms.  Radial cutoffs, followed by local smooth approximation,
are dense in the full-space form norm.  Indeed, for cutoffs \(\chi_R\)
with \(|\nabla\chi_R|\le C/R\), the additional gradient term is bounded by
\(C R^{-2}\int_{|x|>R}w|\xi|^2\).  For odd functions the analogous
approximation on \(\Omega_+\) uses their zero trace on the median plane.

\subsection{The constrained angle}

\begin{proposition}[Existence and regularity of the angle]
\label{prop:angle-minimizer}
Let \(n\) be a positive, continuous, even density with
\(\int n=3\), \(n\in H^2_{\mathrm{loc}}(\R^3)\), and
\(\sqrt n\in H^1(\R^3)\).  The problem
\begin{equation}\label{eq:angle-minimum}
 \begin{aligned}
 &\min_{\vartheta\in H^1(n;\R)}
       \int_{\R^3}n|\nabla\vartheta|^2,\\
 &\text{subject to }\quad
       \vartheta\text{ odd},\qquad
       \int_{\R^3}n\sin^2\vartheta=1
 \end{aligned}
\end{equation}
has a minimizer \(\theta\) satisfying
\(0\le\theta\le\pi/2\) on \(\Omega_+\).
Every such folded minimizer has a multiplier \(d>0\) and satisfies
\begin{equation}\label{eq:angle-euler}
 L_n\theta:=n^{-1}\operatorname{div}(n\nabla\theta)
       =-d\sin(2\theta)
 \quad\hbox{on }\R^3.
\end{equation}
It belongs locally to \(W^{2,p}\) for every \(p<6\), and
\[
 0<\theta<\pi/2\quad\hbox{on }\Omega_+.
\]
In particular, it is locally \(C^{1,\alpha}\) for every
\(0<\alpha<1/2\).  Axial symmetry of \(n\) is not required for this
proposition, and axial symmetry of \(\theta\) is not asserted.
\end{proposition}

\begin{proof}
On the upper half-space the Lipschitz map
\(t\mapsto\arcsin|\sin t|\) preserves \(\sin^2t\), preserves zero
trace, and does not increase gradient energy.  Extend the folded
function oddly.  Thus every competitor can be replaced by one taking
values in \([0,\pi/2]\) above the plane.

The admissible class is nonempty.  Smooth odd angles that equal
\(\pi/2\) on increasingly large compact subsets of \(\Omega_+\)
have constraint values approaching \(3\).  For one such angle the
constraint is greater than \(1\); multiplying it by a parameter in
\([0,1]\) and using continuity produces the required value \(1\).
A folded minimizing sequence is bounded in \(H^1\) on every compact
set, since \(n\) has a positive local lower bound.  Local compactness
gives an odd limit with zero trace on the plane.  The bound
\(\sin^2\theta\le1\) and integrability of \(n\) preserve the constraint
under exhaustion of space.  Lower semicontinuity proves minimality.
The bound on the angle also gives its membership in \(L^2(n)\).

Folding any perturbed competitor shows that a folded minimum is a
minimum for unrestricted odd variations.  The constraint derivative
is nonzero on some compact smooth odd variation.  Otherwise
\(\sin(2\theta)=0\) almost everywhere on \(\Omega_+\).
A locally Sobolev function with values in the discrete set
\(\{0,\pi/2\}\) has zero gradient and is constant on that connected
set; the trace and constraint exclude both possibilities.
The Lagrange-multiplier rule, normalized by a factor of two, therefore
gives
\[
 \int n\nabla\theta\cdot\nabla\varphi
       =d\int n\sin(2\theta)\varphi
\]
for compact odd tests.  The same identity for even tests is zero on
both sides by parity, proving the full-space equation.

Testing with cutoffs times the bounded function \(\theta\) gives
\begin{equation}\label{eq:angle-positive-multiplier}
 \int n|\nabla\theta|^2
       =d\int n\theta\sin(2\theta).
\end{equation}
The cutoff error tends to zero by finite gradient energy and
\(\int n<\infty\).  The left side is positive.  The integral on the
right without \(d\) is positive by folding, oddness, and the preceding
nonconstancy argument.  Hence \(d>0\).

Here are the local estimates needed below.  Positivity of \(n\) and
\(n\in H^2_{\mathrm{loc}}\) imply
\(\nabla\log n\in L^6_{\mathrm{loc}}\).  Rewrite
\Cref{eq:angle-euler} by testing with \(\varphi/n\), where
\(\varphi\) is compactly supported and smooth.  The positive local
lower bound and \(H^2_{\mathrm{loc}}\) regularity of \(n\) make this
a legitimate local \(H^1\) test.  Thus, in distributions,
\[
 \Delta\theta
   =-(\nabla\log n)\cdot\nabla\theta-d\sin(2\theta).
\]
Its right side initially lies in \(L^{3/2}_{\mathrm{loc}}\).
Interior estimates for the ordinary Laplacian on nested compact
sets~\cite[pp.~6--7]{Yuan19} give
\[
 W^{2,3/2}_{\mathrm{loc}}
 \ \Longrightarrow\ \nabla\theta\in L^3_{\mathrm{loc}}
 \ \Longrightarrow\ H^2_{\mathrm{loc}}
 \ \Longrightarrow\ W^{2,3}_{\mathrm{loc}}.
\]
The last space gives \(\nabla\theta\in L^q_{\mathrm{loc}}\) for every
finite \(q\), and a final application of the equation gives
\(W^{2,p}_{\mathrm{loc}}\) for every \(p<6\).
These estimates are uniform when the local \(H^2\) bounds and positive
lower bounds for \(n\), the local angle energy, and \(d\) are uniform.
They apply across the median plane because the equation has already
been extended to full space.

Finally, \(\theta\ge0\) solves an equation with bounded coefficient
\(d\sin(2\theta)/\theta\), and \(\delta=\pi/2-\theta\ge0\) on
\(\Omega_+\) satisfies
\begin{equation}\label{eq:delta-equation}
 L_n\delta=q\delta,\qquad
 q=d\,\frac{\sin(2\delta)}{\delta}\in[0,2d],
\end{equation}
with the continuous value \(2d\) at \(\delta=0\).
On each compact ball the divergence-form coefficient \(nI\) is
bounded and uniformly elliptic, and the zeroth-order coefficients
are bounded.  The local Harnack inequality therefore excludes an
interior zero of either nonnegative function
\cite[Lecture~17, p.~199, and Lecture~18, Theorem~1]{SimonPDE}:
the resulting constant would violate the trace or the constraint.
\end{proof}

The hypotheses of \Cref{prop:angle-minimizer} hold for the densities
under consideration by
\Cref{lem:density-regularity,prop:density-properties}.
In the remainder of the qualitative argument, \(D\) is fixed and
sufficiently large and \(\lambda\) tends to zero.

\begin{proposition}[Convergence of every folded minimum]
\label{prop:angle-convergence}
For every choice of folded minimizing angles for \(n_\lambda\),
\begin{equation}\label{eq:angle-convergence}
 \theta_\lambda\longrightarrow
 \theta_*=\arctan\!\left(\frac{e}{\sqrt2\,b}\right),
 \qquad
 d_\lambda\longrightarrow d_*=\epsilon_e-\epsilon_b.
\end{equation}
The angle convergence is weak in local \(H^1\), strong in local
\(L^2\), and in \(C^1\) on every fixed compact set.  In particular,
all these conclusions hold uniformly over the set of folded
minimizers in the sequential sense.
\end{proposition}

\begin{proof}
At \(\lambda=0\), let
\(g=\sqrt{n_0}\cos\vartheta\) and
\(h=\sqrt{n_0}\sin\vartheta\) for an admissible folded angle.
Their squared norms are \(2\) and \(1\); they are orthogonal by
parity.  The determinant with occupied spin orbitals
\((g/\sqrt2)\uparrow,(g/\sqrt2)\downarrow,h\uparrow\)
has density \(n_0\) and kinetic energy
\[
 \frac12\int\bigl(|\nabla g|^2+|\nabla h|^2\bigr)
 =\frac12\int|\nabla\sqrt{n_0}|^2
       +\frac12\int n_0|\nabla\vartheta|^2.
\]
Adding the fixed integral \(\int u n_0\) and using the simple
three-electron ground state in the indicated spin sector shows that
the minimum is attained uniquely by
\(g=\sqrt2 b,\ h=e\), with the signs fixed by folding.
Indeed, equality fixes the occupied spin-orbital subspace; its
one-dimensional down-spin part fixes \(g\), and the remaining
up-spin orbital then fixes \(h\).
This gives precisely \(\theta_*\).

To obtain a recovery sequence for small \(\lambda\), first multiply
the bounded \(\theta_*\) by radial cutoffs.  Its finite angle energy
and the tail of \(n_0\) make the resulting energy and constraint
errors tend to zero.  Approximate on the remaining compact set by
smooth odd functions.  A fixed compact odd variation with nonzero
constraint derivative corrects the constraint exactly, both at zero
and at nearby \(\lambda\).  Local convergence of \(n_\lambda\)
then gives the energy upper bound for minima.  Folding the corrected
competitor, if necessary, only decreases its energy.

The minimizing angles consequently have uniformly bounded local
\(H^1\) norms.  Common tightness of \(n_\lambda\), supplied by
\Cref{lem:density-regularity}, preserves the constraint for any local
limit.  Lower semicontinuity and the recovery bound make that limit
a minimum at zero coupling, hence \(\theta_*\).
For a fixed compact odd \(\varphi\) with nonzero limiting constraint
derivative, the weak equation gives
\[
 d_\lambda
 =\frac{\int n_\lambda\nabla\theta_\lambda\cdot\nabla\varphi}
        {\int n_\lambda\sin(2\theta_\lambda)\varphi}.
\]
Its numerator and denominator converge.  The limiting orbital
equations for \(b,e\) identify the quotient as \(d_*\).
The uniform local estimates in
\Cref{prop:angle-minimizer}, together with compact Sobolev embedding,
now give \(C^1\) convergence on compact sets.
Applying this argument to every sequence of minimizers proves the
asserted uniformity over all folded choices.
\end{proof}

\begin{proposition}[The two orbital equations]
\label{prop:orbital-equations}
For a folded minimizer define
\begin{equation}\label{eq:angle-orbitals}
 g=\sqrt n\cos\theta,\qquad h=\sqrt n\sin\theta.
\end{equation}
Then \(g,h\in H^1(\R^3)\cap H^2_{\mathrm{loc}}(\R^3)\),
\(g\) is even and strictly positive, and \(h\) is odd and strictly
positive on \(\Omega_+\).  Moreover,
\[
 \|g\|_2^2=2,\qquad \|h\|_2^2=1,\qquad
 \langle g,h\rangle=0,\qquad g^2+h^2=n.
\]
The real function
\begin{equation}\label{eq:formal-potential}
 V=\frac{\Delta g}{2g}
   =\frac{\Delta\sqrt n}{2\sqrt n}
        -\frac{|\nabla\theta|^2}{2}+d\sin^2\theta
\end{equation}
belongs to \(L^2_{\mathrm{loc}}(\R^3)\), and
\begin{equation}\label{eq:orbital-equations}
 (-\Delta/2+V)g=0,\qquad
 (-\Delta/2+V)h=dh
\end{equation}
in distributions.  On the median plane, \(g=\sqrt n\) and \(h=0\).
\end{proposition}

\begin{proof}
The norm, parity, and positivity statements follow from the constraint
and \Cref{prop:angle-minimizer}.  The Sobolev chain rule gives
\begin{equation}\label{eq:angle-kinetic-identity}
 |\nabla g|^2+|\nabla h|^2
       =|\nabla\sqrt n|^2+n|\nabla\theta|^2,
\end{equation}
so both orbitals belong to \(H^1\).  Positivity and local \(H^2\)
regularity of \(n\), together with the angle estimates, give local
\(H^2\) regularity of the orbitals.  Differentiating
\Cref{eq:angle-orbitals} and using \Cref{eq:angle-euler} proves
\Cref{eq:formal-potential,eq:orbital-equations}.
Since the continuous function \(g\) has a positive lower bound on
each compact set, \(\Delta g/(2g)\) is locally square integrable.
\end{proof}

\subsection{A strip comparison and exterior coercivity}

The orbital equations alone do not establish their spectral ordering.
We first control tests supported far from the nuclei, including tests
with nonzero trace on the median plane.  This will prevent weighted
norm from escaping when we compare the even spectrum with its
zero-coupling limit.

Fix throughout this subsection
\[
 L=10,\qquad \mu=\frac1{10}.
\]
The density is axisymmetric, so it may be written as \(n(p,z)\).
This notation will not be imposed on the angle.

\begin{lemma}[The one-dimensional comparison function]
\label{lem:strip-solution}
Suppose, for \(p\ge P_0\) and \(0\le z\le L\), that
\(m(p,z)=e^{2kp}n(p,z)>0\) satisfies
\begin{equation}\label{eq:strip-comparison-data}
 \begin{aligned}
 |m_p|+|m_{pp}|&\le\eta m,\qquad 0<\eta\le1,\\
 \frac{m(p,t)}{m(p,s)}&\le C
       \qquad(0\le t\le s\le L).
 \end{aligned}
\end{equation}
Let \(U(p,\cdot)\) solve
\begin{equation}\label{eq:strip-ode}
 n^{-1}(nU_z)_z=\mu^2U,\qquad U(p,0)=1,\quad U(p,L)=0.
\end{equation}
It is positive and decreasing on \(0\le z<L\), and
\begin{equation}\label{eq:strip-horizontal-derivatives}
 |U_p|+|U_{pp}|\le C_2\eta U,
\end{equation}
where \(C_2\) depends only on \(C,L,\mu\).
When \(C\le100\), one may take \(C_2<10^{30}\).
\end{lemma}

\begin{proof}
Solve backward with final flux equal to one.  If \(J\) is the positive
flux magnitude, the unnormalized solution and its flux obey
\begin{equation}\label{eq:strip-volterra}
 \begin{aligned}
 \widetilde U(p,z)&=\int_z^L\frac{J(p,s)}{m(p,s)}\,ds,\\
 J(p,z)&=1+\mu^2\int_z^L\int_t^L
                \frac{m(p,t)}{m(p,s)}J(p,s)\,ds\,dt.
 \end{aligned}
\end{equation}
The successive-integral series is positive.  Its term of order \(l\)
is bounded by
\[
 \frac{(\mu^2 C)^l(L-z)^{2l}}{(2l)!}.
\]
It therefore converges uniformly, defines \(J\ge1\), and gives the
unique solution after normalization
\(U(p,z)=\widetilde U(p,z)/\widetilde U(p,0)\).
The formulas give \(U>0\) and \(U_z<0\).

For a kernel \(K=m(p,t)/m(p,s)\),
\[
 |K_p|\le2\eta K,\qquad |K_{pp}|\le6\eta K.
\]
The product of \(l\) such kernels has derivatives of either positive
order at most two bounded by \(20(1+l)^2\eta\) times that product.
Termwise differentiation of the series is consequently justified.
For \(C\le100\), the undifferentiated terms are bounded by
\(10^{2l}/(2l)!\), and
\[
 \sum_{l\ge0}20(1+l)^2\frac{10^{2l}}{(2l)!}<10^8.
\]
Since \(J\ge1\), these estimates bound its first two \(p\)-derivatives
relative to \(J\).  Also
\(|(m^{-1})_p|\le\eta m^{-1}\) and
\(|(m^{-1})_{pp}|\le3\eta m^{-1}\).
Integrating the positive expression for \(\widetilde U\), and then
differentiating its quotient by \(\widetilde U(p,0)\), proves
\Cref{eq:strip-horizontal-derivatives}; the displayed numerical bounds
give \(C_2<10^{30}\) with ample slack.
For general fixed \(C\) the same convergent series gives a finite
constant.  Only the one-sided ratios in
\Cref{eq:strip-comparison-data} were used.
\end{proof}

\Needspace{5\baselineskip}
\begin{proposition}[Uniform exterior gap]\label{prop:exterior-gap}
For sufficiently large \(D\), there are \(\lambda_0>0\) and \(R_0<\infty\)
such that every folded minimizing angle for
\(0\le\lambda<\lambda_0\) satisfies
\begin{equation}\label{eq:exterior-weighted-gap}
 \frac12\int_{\Omega_+}g^2|\nabla\xi|^2
       \ge0.03\int_{\Omega_+}g^2|\xi|^2
\end{equation}
for tests supported in \(\{|x|>R_0\}\).
No zero trace on \(z=0\) is required.
The inequality holds throughout the corresponding weighted form domain.
\end{proposition}

\begin{proof}
Set \(\delta=\pi/2-\theta\) on \(\Omega_+\).
By \Cref{eq:delta-equation},
\[
 L_n\delta=q\delta,\quad 0\le q\le2d,\qquad
 \frac4{\pi^2}n\delta^2\le g^2=n\sin^2\delta\le n\delta^2.
\]
Choose \(D\) large enough that \(d_*\) and \(\gamma-k_0\) are small.
The multiplier convergence and the density estimates
\Cref{eq:density-drift,eq:strip-data} allow us to choose the small
constants below uniformly for sufficiently small \(\lambda\).

First, \(U\) from \Cref{lem:strip-solution} is a subsolution in the
strip.  In cylindrical coordinates,
\[
 L_nU=\mu^2U+U_{pp}
          +(p^{-1}+\partial_p\log n)U_p.
\]
The horizontal coefficient is bounded sufficiently far out by
\Cref{eq:density-drift}.  Thus
\Cref{eq:strip-horizontal-derivatives}, with \(\eta\) small, and
\(q\le2d\), with \(d\) small, imply
\begin{equation}\label{eq:strip-subsolution}
 (L_n-q)U\ge0.
\end{equation}

We next supply an inner boundary for comparison.  At zero coupling
put \(\delta_*=\pi/2-\theta_*\).  On \(z>0\),
\[
 \delta_*=\arctan(\sqrt2 b/e).
\]
The scalar exterior estimates in
\Cref{eq:unperturbed-strip-ratio} show that, uniformly for
\(0<z\le L\), the first two horizontal logarithmic derivatives of
\(b/e\) are small when \(D\) and then the starting radius are large.
This remains true down to \(z=0\): the simple odd factor of \(e\)
is canceled before taking a horizontal derivative.
If \(t>0\), the ratios of
\(t(\arctan t)'\) and \(t^2(\arctan t)''\) to \(\arctan t\)
are bounded.  Consequently
\[
 \frac{|(\delta_*)_p|+|(\delta_*)_{pp}|}{\delta_*}
\]
is uniformly small in the strip.  Combining this with its angle
equation gives, for a fixed sufficiently large \(P_1\),
\begin{equation}\label{eq:inner-strip-inequality}
 n_0^{-1}\partial_z(n_0\partial_z\delta_*)
       <\mu^2\delta_*
       \quad(p=P_1,\ 0\le z\le L).
\end{equation}
The strict inequality follows by making the horizontal error and
\(2d_*\) smaller than \(\mu^2\).
One-dimensional comparison with \(U_0(P_1,\cdot)\) now gives
\(\delta_*/(\pi/2)>U_0\) for \(0<z\le L\), and the derivative of
their difference at \(z=0\) is strictly positive.
Its quotient by \(z\) therefore has a positive minimum on
\(0<z\le L\).  By \Cref{prop:angle-convergence} and continuity of the
one-dimensional problem in its positive coefficient, this boundary
ordering persists for small \(\lambda\), uniformly over all folded
minimizers and all azimuths.
For this use of continuity, \Cref{eq:strip-volterra} gives uniform
convergence of \(U_\lambda\) and \((U_\lambda)_z\) on the fixed segment
from uniform convergence and a positive lower bound for its density.

It follows that
\begin{equation}\label{eq:strip-angle-comparison}
 \frac{\delta(x)}{\pi/2}\ge U(p,z)
       \quad(p\ge P_1,\ 0\le z\le L).
\end{equation}
Here is a comparison proof that also addresses infinity.
The difference \(v=\delta/(\pi/2)-U\) satisfies
\((L_n-q)v\le0\) and has nonnegative finite-boundary values.
Test this inequality with a horizontal cutoff squared times its
negative part, with the appropriate negative sign.
The resulting Caccioppoli inequality is
\[
 \int n\chi_R^2|\nabla v_-|^2
       +\int nq\chi_R^2v_-^2
       \le4\int n v_-^2|\nabla\chi_R|^2.
\]
The right side tends to zero because \(v_-\) is bounded and \(n\)
is integrable.  Thus \(v_-\) is constant and its zero boundary trace
makes it zero.  No value at infinity has been imposed.
This proof also makes clear that \(\theta\) need not be axisymmetric:
the coefficient \(q\) may depend on azimuth.
Differentiating the boundary inequality at \(z=0\) gives
\begin{equation}\label{eq:strip-flux-bound}
 -\frac{\delta_z}{\delta}\le-U_z(p,0).
\end{equation}

For a real smooth test \(\xi\) compactly supported sufficiently far
out, put \(F_1=\delta\xi\).  Integration by parts, with outward
normal \(-e_z\) on the plane, yields
\begin{equation}\label{eq:delta-transform-boundary}
 \int_{\Omega_+}n\delta^2|\nabla\xi|^2
 =\int_{\Omega_+}n\bigl(|\nabla F_1|^2+q|F_1|^2\bigr)
  -\int_{z=0}n\left(-\frac{\delta_z}{\delta}\right)|F_1|^2.
\end{equation}
The energy of the solution of \Cref{eq:strip-ode} is
\[
 \int_0^L n(U_z^2+\mu^2U^2)\,dz=-n(p,0)U_z(p,0).
\]
It minimizes this energy among functions with the same endpoint
values.  Use the trial function
\(\cos(az)F_1\), \(a=\pi/(2L)\), for each fixed horizontal point.
The exact identity
\[
 |F_1'|^2+a^2|F_1|^2
       -|(\cos(az)F_1)'|^2
       =|\sin(az)F_1'+a\cos(az)F_1|^2
\]
gives
\begin{equation}\label{eq:strip-boundary-cost}
 n(p,0)(-U_z(p,0))|F_1(0)|^2
 \le\int_0^L n\bigl(
        |\partial_zF_1|^2+(\mu^2+a^2)|F_1|^2\bigr)\,dz.
\end{equation}
Thus the boundary term spends at most the vertical gradient energy
in the strip, together with the indicated \(L^2\) cost.

The remaining gradient energy controls
\(\int_{\Omega_+}n|Y\cdot\nabla F_1|^2\), where \(Y\) is the
horizontal radial unit vector for \(z\le p/4\), the spatial radial
unit vector for \(z\ge p/2\), and a smooth convex interpolation
between them.  More explicitly, choose \(\zeta(t)\in[0,1]\), equal to
zero for \(t\le1/4\) and one for \(t\ge1/2\), with
\(|\zeta'|\le8\), and set
\[
 Y=(1-\zeta(z/p))\frac{(x_1,x_2,0)}p
            +\zeta(z/p)\frac{x}{r}.
\]
On the axis use \(Y=x/r\).  The field has norm at most one and is
tangential to the median plane.  On the interpolation region,
\(|\nabla(z/p)|\le5/(4r)\) and \(1/p\le\sqrt5/(2r)\).
The product rule therefore gives
\(\lvert\operatorname{div}Y\rvert\le100/r\).
For sufficiently distant support it is horizontal throughout
\(0\le z\le L\).  In the interpolation region \(p/r\ge2/\sqrt5\);
the two drift bounds in \Cref{eq:density-drift}, with \(k\) close
to one, therefore imply
\begin{equation}\label{eq:outward-weighted-divergence}
 n^{-1}\operatorname{div}(nY)\le-1
\end{equation}
after increasing the fixed starting radius.
There is no boundary contribution when integrating
\(\operatorname{div}(nY|F_1|^2)\).  Cauchy--Schwarz gives
\[
 \int_{\Omega_+}n|Y\cdot\nabla F_1|^2
       \ge\frac14\int_{\Omega_+}n|F_1|^2.
\]
Dropping the nonnegative \(q\) term in
\Cref{eq:delta-transform-boundary} and using
\Cref{eq:strip-flux-bound,eq:strip-boundary-cost} proves
\[
 \int_{\Omega_+}n\delta^2|\nabla\xi|^2
 \ge\left(\frac14-\mu^2-\frac{\pi^2}{4L^2}\right)
             \int_{\Omega_+}n\delta^2|\xi|^2.
\]
The comparison of \(g^2\) with \(n\delta^2\) yields the coefficient
\begin{equation}\label{eq:exterior-gap-constant}
 \frac2{\pi^2}
 \left(\frac14-\frac1{100}-\frac{\pi^2}{400}\right)
       >0.0436>0.03.
\end{equation}
This proves \Cref{eq:exterior-weighted-gap}.  Real and imaginary
parts give complex tests; the form-domain approximation described
at the beginning of the section gives the general statement.
\end{proof}

\subsection{The full weighted spectral gap}

\begin{proposition}[The gap at the second orbital]
\label{prop:weighted-gap}
For sufficiently large \(D\) and all sufficiently small
\(\lambda\ge0\), every folded minimizing angle satisfies
\begin{equation}\label{eq:weighted-gap}
 \frac12\int_{\R^3}g^2|\nabla\xi|^2
 \ge d\left(
      \int_{\R^3}g^2|\xi|^2
       -\frac{\left|\int_{\R^3}g^2\xi\right|^2}
                    {\int_{\R^3}g^2}\right),
 \qquad \xi\in H^1(g^2).
\end{equation}
\end{proposition}

\begin{proof}
Put \(w=g^2\).  On the upper half-space, \(f=h/g>0\) satisfies
\[
 -\frac1{2w}\operatorname{div}(w\nabla f)=df.
\]
For compact tests inside \(\Omega_+\), the ground-state
transform~\cite[Section~2.3, equation~(2.12)]{FS08} gives
\begin{equation}\label{eq:odd-ground-transform}
 \frac12\int_{\Omega_+}w|\nabla\xi|^2
       -d\int_{\Omega_+}w|\xi|^2
 =\frac12\int_{\Omega_+}wf^2
                 |\nabla(\xi/f)|^2\ge0.
\end{equation}
Odd tests have zero trace at the plane.  The left side of
\Cref{eq:odd-ground-transform} is continuous in the weighted form
norm.  Cutoff and local zero-trace approximation therefore extend
its nonnegativity to all odd
functions in \(H^1(w)\), proving the odd part of the desired gap.

We prove the even part by controlling possible loss of weighted norm
at infinity.  Let \(w_0=2b^2\).  The positive-orbital transform for
\(b\) identifies \(\mathcal E_{w_0}\) with the quadratic form of
\(h_D-\epsilon_b\), under \(\xi\mapsto\sqrt2b\,\xi\).
The identity first holds for compact smooth tests.  Their density in
\(H^1(w_0)\), together with closedness of the Schr\"odinger form,
extends it to the full weighted domain: the transformed approximants
are Cauchy in the form norm, and their \(L^2\) limit is
\(\sqrt2b\,\xi\).  No global bound on \(\nabla\log b\) is needed.
By \Cref{prop:two-centre-spectrum}, there is an even spectral gap
\(\Delta_{\mathrm{ev}}>0.2\):
\begin{equation}\label{eq:unperturbed-even-gap}
 \mathcal E_{w_0}(\xi)\ge
       \Delta_{\mathrm{ev}}\int w_0|\xi|^2
       \quad\hbox{if \(\xi\) is even and }\int w_0\xi=0.
\end{equation}
In particular it is strictly greater than \(d_*\), and we choose
\(D\) so that \(d_*<0.03\).

If the even assertion failed arbitrarily close to zero coupling,
there would be \(\lambda_j\downarrow0\), folded minimizers, and real
even tests with
\[
 \int w_j\xi_j=0,\qquad
 \int w_j|\xi_j|^2=1,\qquad
 \limsup_j\mathcal E_{w_j}(\xi_j)\le d_*,
 \quad w_j=g_{\lambda_j}^2.
\]
Local uniform convergence to \(w_0\), and its positive local lower
bounds, give a subsequence converging weakly in local \(H^1\)
and strongly in local \(L^2\) to an even \(\xi\).
Write \(s=\int w_0|\xi|^2\le1\).  The limit has mean zero, since
\[
 \left|\int_{|x|>R}w_j\xi_j\right|
       \le\left(\int_{|x|>R}w_j\right)^{1/2},
\]
and the tails on the right are uniformly small:
\(w_j\le n_{\lambda_j}\), and the latter densities are tight.
No tightness of \(w_j|\xi_j|^2\) is assumed.

Choose a smooth radial quadratic partition
\(\chi_R^2+\rho_R^2=1\), with \(\chi_R=1\) on \(r\le R\),
\(\chi_R=0\) on \(r\ge2R\), and derivatives of size \(O(R^{-1})\).
Expanding gradients gives the weighted localization identity, as in
the usual IMS formula~\cite[Section~3, Lemma~3.1]{Simon83}:
\[
 \mathcal E_{w_j}(\xi_j)
 =\mathcal E_{w_j}(\chi_R\xi_j)
  +\mathcal E_{w_j}(\rho_R\xi_j)
  -\frac12\int w_j
       (|\nabla\chi_R|^2+|\nabla\rho_R|^2)|\xi_j|^2.
\]
The last term is \(O(R^{-2})\).  For \(R\ge R_0\),
\Cref{prop:exterior-gap}, applied to both half-spaces, bounds
the second term below by \(0.03\|\rho_R\xi_j\|_{L^2(w_j)}^2\).
Local lower semicontinuity and local strong convergence therefore
give
\[
 \liminf_j\mathcal E_{w_j}(\xi_j)
 \ge \mathcal E_{w_0}(\chi_R\xi)
      +0.03\left(1-\int w_0|\chi_R\xi|^2\right)-O(R^{-2}).
\]
Local lower semicontinuity also gives finite global energy to the
limit.  Letting \(R\to\infty\) and using
\Cref{eq:unperturbed-even-gap}, we obtain
\[
 d_*\ge\mathcal E_{w_0}(\xi)+0.03(1-s)
       \ge\Delta_{\mathrm{ev}}s+0.03(1-s)>d_*,
\]
a contradiction for every \(s\in[0,1]\).
This proves the even assertion uniformly for all folded minimizers.
Finally, reflection separates the even and odd parts orthogonally
in the weighted norm and energy.  The odd part has zero weighted
mean.  Combining the two inequalities, and then the real and
imaginary parts, proves \Cref{eq:weighted-gap}.
\end{proof}

\subsection{A finite collection of sufficient tests}

The following version will be used for the exact charge prescription.
It avoids requiring that the selected coupling belong to an
unspecified perturbation interval.

\begin{lemma}[A quantitative local test for the gap]
\label{lem:quantitative-gap}
Let \(b,e,n_0,\theta_*\) be the zero-coupling data for a fixed \(D\)
whose even gap satisfies \(\Delta_{\mathrm{ev}}>0.2\).
Let \(n\) satisfy the hypotheses of \Cref{prop:angle-minimizer},
be axisymmetric, and let \(\theta\) be a folded minimizer.  Set
\[
 \eta_0=10^{-60},\qquad M\ge10^4,\qquad S=100M.
\]
Suppose \(n\) is \(C^2\) on a neighborhood of \(\{r\ge M\}\),
\(0.99<k<1.01\), \(0<d<0.001\), and the following tests hold.
\begin{enumerate}[label=\textup{(\roman*)}]
\item The two drift bounds of \Cref{eq:density-drift} hold with
\(\eta=\eta_0\) for \(r\ge M\) on their stated domains.
For \(p\ge M,\ 0\le z\le L\), the bounds
\Cref{eq:strip-comparison-data} hold with
\(\eta=\eta_0\) and \(C=100\).
\item Let \(U,U_0\) be the strip solutions for \(n,n_0\), and suppose
\[
 \alpha=\inf_{0<z\le L}
     \frac{(2/\pi)\delta_*(M,z)-U_0(M,z)}{z}>0,
 \qquad \delta_*=\pi/2-\theta_*.
\]
Uniformly on the cylinder \(p=M,\ 0\le z\le L\), including all
azimuths for \(\theta\), one has
\[
 |\partial_z(\theta-\theta_*)|
       +|\partial_z(U-U_0)|<\alpha/2.
\]
\item On \(r\le2S\), the local weight comparison is
\[
 \frac12\le\frac{g^2}{2b^2}\le2.
\]
\end{enumerate}
Then \Cref{eq:exterior-weighted-gap} holds for tests supported in
\(r>S\), and the full gap \Cref{eq:weighted-gap} holds.
\end{lemma}

\begin{proof}
The strip derivative constant in \Cref{lem:strip-solution} is less
than \(10^{30}\).  The drift coefficient in the strip has absolute
value less than \(3\), so the horizontal error in \(L_nU\) is at
most \(4\cdot10^{30}\eta_0U\).  Since \(q\le2d<0.002\) and
\(\mu^2=0.01\), this proves \Cref{eq:strip-subsolution}
for every \(p\ge M\).
At the inner cylinder the function
\[
 \frac2\pi\delta-U
       -\left(\frac2\pi\delta_*-U_0\right)
\]
vanishes at \(z=0\) and has \(z\)-derivative of absolute value less
than \(\alpha/2\).  Its reference term is at least \(\alpha z\).
Thus the required inner boundary ordering is strict.
The same cutoff comparison proves
\Cref{eq:strip-angle-comparison} for \(p\ge M\).

For \(r\ge S\), the vector field used above is horizontal throughout
the strip.  Its weighted divergence is bounded by
\[
 -2k\,\frac2{\sqrt5}+\eta_0+\frac{100}{r}<-1.
\]
All boundary support lies at \(p>M\).  The flux calculation and
\Cref{eq:exterior-gap-constant} consequently give the exterior
constant \(0.03\) from radius \(S\).

It remains to prove the even part of the full gap without a
perturbative compactness argument.  Choose a nonnegative smooth
radial quadratic partition \(\chi^2+\rho^2=1\) between \(S\) and
\(2S\), with
\[
 \frac12\bigl(|\nabla\chi|^2+|\nabla\rho|^2\bigr)
       \le\frac{10}{S^2}.
\]
For an even test \(\xi\in H^1(g^2)\), put \(w=g^2\), \(w_0=2b^2\),
and choose
\[
 c=\frac{\int w_0\chi\xi}{\int w_0\chi},
 \qquad f=\xi-c.
\]
The denominator is positive and \(\int w_0\chi f=0\).
By the local weight comparison and
\Cref{eq:unperturbed-even-gap},
\[
 \begin{aligned}
 \mathcal E_w(\chi f)
 &\ge\frac12\mathcal E_{w_0}(\chi f)
 \ge\frac{\Delta_{\mathrm{ev}}}{2}
               \|\chi f\|_{L^2(w_0)}^2\\
 &\ge\frac{\Delta_{\mathrm{ev}}}{4}
               \|\chi f\|_{L^2(w)}^2
 \ge0.05\|\chi f\|_{L^2(w)}^2.
 \end{aligned}
\]
The exterior estimate gives
\(\mathcal E_w(\rho f)\ge0.03\|\rho f\|_{L^2(w)}^2\).
Although \(f\) need not have compact support, it belongs to
\(H^1(w)\), because \(\int w=2\); thus the established cutoff
extension applies.  Localization yields
\[
 \mathcal E_w(\xi)=\mathcal E_w(f)
 \ge\left(0.03-\frac{10}{S^2}\right)\|f\|_{L^2(w)}^2
 \ge d\,\Var_w(\xi).
\]
The last step uses \(d<0.001\), the stated lower bound on \(S\),
and the fact that subtracting any constant gives squared norm at
least the weighted variance.  The odd ground-state transform
\Cref{eq:odd-ground-transform} does not require small coupling and
completes the proof, including complex tests.
\end{proof}
